% Options for packages loaded elsewhere
\PassOptionsToPackage{unicode}{hyperref}
\PassOptionsToPackage{hyphens}{url}
\PassOptionsToPackage{dvipsnames,svgnames,x11names}{xcolor}
\documentclass[
  11pt,
  letterpaper,
]{article}
\usepackage{xcolor}
\usepackage[margin=1in]{geometry}
\usepackage{amsmath,amssymb}
\setcounter{secnumdepth}{-\maxdimen} % remove section numbering
% Classic Computer Modern; no system-font or Unicode-math fallback.
\usepackage{iftex}
\usepackage[OT1]{fontenc}
\renewcommand{\rmdefault}{cmr}
\renewcommand{\sfdefault}{cmss}
\renewcommand{\ttdefault}{cmtt}

\ifPDFTeX\else
  % xetex/luatex font selection
\fi
% Use upquote if available, for straight quotes in verbatim environments
\IfFileExists{upquote.sty}{\usepackage{upquote}}{}
\IfFileExists{microtype.sty}{% use microtype if available
  \usepackage[]{microtype}
  \UseMicrotypeSet[protrusion]{basicmath} % disable protrusion for tt fonts
}{}
\makeatletter
\@ifundefined{KOMAClassName}{% if non-KOMA class
  \IfFileExists{parskip.sty}{%
    \usepackage{parskip}
  }{% else
    \setlength{\parindent}{0pt}
    \setlength{\parskip}{6pt plus 2pt minus 1pt}}
}{% if KOMA class
  \KOMAoptions{parskip=half}}
\makeatother
\setlength{\emergencystretch}{3em} % prevent overfull lines
\providecommand{\tightlist}{%
  \setlength{\itemsep}{0pt}\setlength{\parskip}{0pt}}
\usepackage{microtype}
\usepackage{amssymb}
\setlength{\emergencystretch}{3em}
\allowdisplaybreaks[1]
\usepackage{bookmark}
\IfFileExists{xurl.sty}{\usepackage{xurl}}{} % add URL line breaks if available
\urlstyle{same}
% fallback for those not using the hyperref driver hyperxmp:
\makeatletter
\@ifundefined{xmpquote}{\newcommand{\xmpquote}[1]{#1}}{}
\makeatother
\hypersetup{
  pdftitle={Successor truthfulness and particular-code knowledge},
  colorlinks=true,
  linkcolor={Maroon},
  filecolor={Maroon},
  citecolor={Blue},
  urlcolor={blue},
  pdfcreator={LaTeX via pandoc}}

\title{Successor truthfulness and particular-code knowledge}
\usepackage{etoolbox}
\makeatletter
\providecommand{\subtitle}[1]{% add subtitle to \maketitle
  \apptocmd{\@title}{\par {\large #1 \par}}{}{}
}
\makeatother
\subtitle{A negative answer to Alexander\textquotesingle s Open Problem
3.2.1}
\author{}
\date{20 October 2026}

\begin{document}
\maketitle
\begin{abstract}
Alexander's Open Problem 3.2.1 asks whether there can be an infinite
sequence of knowing machines, each knowing the truthfulness and a code
of its successor, without the ordinal knowledge assumed in his preceding
impossibility theorem. We give a negative answer in his recursively
named language and standard canonical semantics. The key observation is
effective: particular-code knowledge for every unary modal formula makes
a single enumeration of the first machine's consequences sufficient to
enumerate the consequences of every later machine, uniformly. The
selected indices retain their predecessors' proofs of correctness.
Applying this observation to a fixed partial truth predicate produces a
uniformly enumerable sequence of sound arithmetic theories, each proving
uniform reflection for its successor. A theorem of Pakhomov and Walsh
excludes that sequence. We give the arithmetic conversion and the
comparison of proof presentations explicitly. The argument requires
neither an ordinal predicate nor an internally proved totality assertion
for the index-selection procedure.
\end{abstract}

\section{1. The question and the
argument}\label{the-question-and-the-argument}

In his 2013 dissertation \emph{The Theory of Several Knowing Machines},
Alexander studies machines that can reason about the knowledge of other
machines. A machine may know that another machine is truthful, that its
knowledge is recursively enumerable, or that a particular program
enumerates its knowledge. These are different requirements.

Chapter 3 proves that an infinite sequence of machines cannot each know
its successor's truthfulness and code when the machines also know
rudimentary facts about recursive ordinals. Section 3.2 asks whether the
ordinal hypothesis can be removed. The precise question is \textbf{Open
Problem 3.2.1, printed page 89}; the definitions of successor
truthfulness and code knowledge are in Definition 3.1.5, printed page 86
{[}A13{]}. The question occurs on page 99 of the 103-page OhioLINK PDF.
It is different from the dissertation's Questions 3-1-1, 3-1-2 and 3-1-3
about well-founded hierarchies.

We answer the question by passing from modal code knowledge to
arithmetic reflection. The proof has three steps.

First, a code for a machine's answers can also be requested for formulas
saying that the machine endorses a code of its successor. Iterating this
request allows us to reach any finite distance from the first machine.
Only the first machine's enumeration is supplied to the procedure.

Second, we retain the arithmetic information needed for a reflection
theorem. For each machine we form a theory consisting of a fixed weak
arithmetic base and all of that machine's designated arithmetical
\(\Pi_3\) consequences. Known successor truthfulness and a known
particular code imply that the preceding machine proves reflection for
this theory.

Third, we arrange the theories into one recursively enumerable
presentation and verify that the reflection statements refer to that
very presentation. The resulting sequence contradicts Pakhomov--Walsh's
arithmetic non-descent theorem {[}PW, Theorem 3.4{]}.

The distinction between an index that happens to be correct and an index
whose correctness is proved by the predecessor is essential in the
second step. The distinction between equal theorem sets and provably
interchangeable proof presentations is essential in the third.

\subsection{Theorem 1. No infinite successor
chain}\label{theorem-1.-no-infinite-successor-chain}

There is no family answering Alexander's Open Problem 3.2.1
affirmatively.

More explicitly, the following weaker collection of requirements is
already impossible. Let \(L\) be the recursive language of arithmetic
with effectively named operators \(K_0,K_1,\ldots\). Let
\(T_0,T_1,\ldots\) be sets of \(L\)-sentences with a common standard
canonical interpretation \(M\). Assume:

\begin{enumerate}
\def\labelenumi{\arabic{enumi}.}
\tightlist
\item
  Every consequence of every \(T_j\) is true in \(M\).
\item
  Every \(T_j\) proves ordinary arithmetic \(\mathsf{PA}\), and the
  consequence set of \(T_0\) is recursively enumerable.
\item
  For every \(j\geq 1\) and every \(L\)-formula \(\varphi(x)\) with
  exactly the free variable \(x\), there is a standard number \(e\) such
  that \[
  T_{j-1}\vdash
  \forall x\bigl(K_j\varphi(x)\leftrightarrow x\in W_{\bar e}\bigr).
  \tag{1}
  \]
\item
  For one fixed arithmetical \(\Pi_3\) partial truth predicate
  \(\tau(x)\), specified in Section 4, every \(T_j\) proves \[
  \forall x\bigl(K_{j+1}\tau(x)\rightarrow\tau(x)\bigr).
  \tag{2}
  \]
\end{enumerate}

Here and below \(\bar e\) is the numeral for \(e\), and \(W_e\) is the
\(e\)-th recursively enumerable set in a fixed effective enumeration.

The theorem does not assume a uniform enumeration of the family. It does
not even need the later consequence sets to be supplied as recursively
enumerable sets: their uniform enumerability will follow from the code
hypothesis. Ordinary arithmetic induction is enough; the source question
assumes the stronger Peano axioms for the full modal language.

\section{2. Canonical knowledge and particular
codes}\label{canonical-knowledge-and-particular-codes}

\subsection{2.1. The consequence
relation}\label{the-consequence-relation}

We use Alexander's consequence relation for his modal language, not a
stronger modal proof calculus. Write \(T\vdash\theta\) when \(\theta\)
follows from \(T\) in that relation. Alexander defines it semantically
and proves compactness and effective enumerability of validities {[}A13,
Theorem 1.2.6{]}. Thus the consequences of a recursively enumerable set
of axioms are recursively enumerable. Ordinary first-order arithmetic
deductions remain valid deductions in this logic.

The canonical interpretation on standard arithmetic is \[
M\models K_j\varphi[s]
\quad\Longleftrightarrow\quad
T_j\vdash\varphi^s,
\tag{3}
\] where \(\varphi^s\) replaces the free variables by the numerals given
by \(s\). This is the interpretation in Alexander's Definitions 1.2.15
and 3.1.4. His Lemma 1.2.12 makes numerical substitution commute with
modal formation in the required way.

Self-doubting knowledge still includes truthfulness:
\(K_j\theta\rightarrow\theta\) holds in the intended structure.
Consequently a \(T_j\)-consequence is true in that structure, by (3).
The adjective \emph{self-doubting} means that the machine need not know
its own truthfulness. Our proof never adds that stronger requirement.
This setting follows Alexander's modification of the earlier
knowing-machine framework of Carlson, Reinhardt and Shapiro; Carlson's
construction is a predecessor, not a theorem proved here {[}C00; A13,
Sections 1.2 and 1.4{]}.

\subsection{2.2. The order of the
quantifiers}\label{the-order-of-the-quantifiers}

For a unary formula \(\varphi(x)\), introduce the unary formula \[
\operatorname{Code}_{j,\varphi}(z)
:=\forall x\bigl(K_j\varphi(x)\leftrightarrow x\in W_z\bigr),
\tag{4}
\] with \(z\) fresh. Bound variables are renamed effectively. When this
formula is itself used as the input of another code request, its free
variable is renamed to the designated input variable.

Alexander's code condition is the truth in \(M\) of \[
\exists e\;K_{j-1}\operatorname{Code}_{j,\varphi}(e).
\tag{5}
\] Because arithmetic is standard, (5) supplies an actual natural number
\(e\); clause (3) then gives precisely (1). It is not the condition \[
K_{j-1}\exists e\;\operatorname{Code}_{j,\varphi}(e).
\tag{6}
\] There is no general rule extracting a particular numeral theorem from
(6). The proof uses (5), which is the condition printed in Definition
3.1.5.

The free variable in (4) is \(z\), even though \(j\) names a modal
operator and \(\varphi\) may be modal. Thus (4), after variable
renaming, is another allowed input to the full unary-formula code
schema. This simple closure property drives the argument.

\section{3. Obtaining all required indices from the first
machine}\label{obtaining-all-required-indices-from-the-first-machine}

Before the general construction, consider reaching \(T_2\). To find a
\(T_1\)-endorsed index for \(T_2\)'s \(\varphi\)-answers, first ask the
root enumeration for a particular code of the following \(T_1\)-column:
\[
\{e:T_1\vdash\operatorname{Code}_{2,\varphi}(\bar e)\}.
\tag{7}
\] This column is nonempty by (1). Enumerate it and take one output.
That output is not just an externally correct index: it occurs because
\(T_1\) proves its correctness. Repeating the operation reaches \(T_3\),
and then every later machine.

\subsection{Lemma 2. Uniform selection with predecessor
endorsement}\label{lemma-2.-uniform-selection-with-predecessor-endorsement}

Fix an enumeration index \(r\) for the consequences of \(T_0\). Under
(1), (3) and soundness, there is a partial recursive procedure \(B\),
obtained effectively from \(r\), such that for every \(j\geq1\) and
every unary modal \(\varphi\), \(B(j,\ulcorner\varphi\urcorner)\) halts
with an output \(e\) satisfying both \[
T_{j-1}\vdash\operatorname{Code}_{j,\varphi}(\bar e)
\tag{8}
\] and \[
W_e=\{m:T_j\vdash\varphi(\bar m)\}.
\tag{9}
\]

\textbf{Proof.} Use deterministic enumerations and a fixed first-output
convention. Define the procedure by recursion on the positive integer
\(j\).

For \(j=1\), enumerate the consequences of \(T_0\) using \(r\). Search
for a sentence of the exact canonical form
\(\operatorname{Code}_{1,\varphi}(\bar e)\), decode its numeral, and
return \(e\).

For the step from \(j\) to \(j+1\), form \[
\chi(z)=\operatorname{Code}_{j+1,\varphi}(z).
\] Compute \(d=B(j,\ulcorner\chi\urcorner)\). Enumerate \(W_d\), and
return its first output \(e\). The recursion decreases the machine
index, although the formula grows.

We prove termination, (8) and (9) together by induction on \(j\),
simultaneously for every unary modal formula. At the base, (1) promises
a sentence for which the root search succeeds. That sentence is a root
theorem, giving (8). Its truth in \(M\), together with standard
arithmetic and (3), gives (9).

For the induction step, \(\chi\) is unary, so the induction hypothesis
gives a terminating computation of \(d\) and \[
W_d=\{e:T_j\vdash\operatorname{Code}_{j+1,\varphi}(\bar e)\}.
\tag{10}
\] Condition (1) says that this set is nonempty. Its enumeration
therefore produces an \(e\). Equation (10) gives (8) for that very
output. Soundness of \(T_j\) then turns the selected theorem into (9),
using (3).

All formula construction, bound-variable renaming, parsing and numeral
decoding are effective. A canonical spelling can be used because the
permitted renamings preserve the relevant formulas. No enumeration of a
later \(T_j\) is supplied to the procedure: later enumerations are
obtained from indices selected by the procedure itself. This proves the
lemma. \hfill\(\square\)

The lemma asserts external termination on the promised inputs. It does
not assert \[
T_0\vdash\forall j\;B(j,\ulcorner\varphi\urcorner)\mathord\downarrow,
\] or any corresponding assertion in \(\mathsf{PA}\). No such theorem
will be required. Notice also that the predecessor endorsement (8) was
proved separately from the extensional equality (9).

\section{4. Arithmetic theories with fixed proof
presentations}\label{arithmetic-theories-with-fixed-proof-presentations}

The next step converts one endorsed arithmetic column into reflection.
We give the details because changing an enumeration or a proof predicate
without justification would leave a gap in the final application of
non-descent.

\subsection{4.1. Arithmetic syntax and partial
truth}\label{arithmetic-syntax-and-partial-truth}

Initially work in the arithmetic language \(L_0=\{0,S,+,\cdot,=\}\).
Designate a recursive class \(\mathcal P_3\) of closed prenex \(\Pi_3\)
formulas with bounded matrices; empty blocks and effective padding are
allowed. A typical member has the form \[
\forall\vec u\;\exists\vec v\;\forall\vec w\;
\delta(\vec u,\vec v,\vec w),
\] where \(\delta\) is bounded. Finite tuples and syntax are coded
arithmetically. Primitive recursive operations used to describe the
coding are abbreviations with their usual arithmetic representations,
not new unproved axioms.

Fix a \(\Pi_3\) partial truth predicate \(\tau(z)\) for this class,
taking the value false on codes outside \(\mathcal P_3\). In particular,
for each fixed \(\sigma\in\mathcal P_3\), \[
\mathsf I\Sigma_1\vdash
\sigma\leftrightarrow\tau(\overline{\ulcorner\sigma\urcorner}).
\tag{11}
\] For each fixed formula this disquotation is verified already in
\(\mathsf I\Sigma_1\): check its bounded matrix and then its three fixed
quantifier blocks. The stronger uniform assertion about truth of all
base axioms below uses \(\mathsf{PA}\). The construction of the
predicate and those uniform equations are recalled in Lemma 3. The basic
use of partial truth and witnessed proofs follows the conventions of
{[}PW, Section 2.1{]}.

Present \(\mathsf I\Sigma_1\) by a fixed recursive set of axioms, each
normalized to a member of \(\mathcal P_3\). For a \(\Sigma_1\) formula
code \(c\), let \(J(c)\) be the normalization of its universally closed
induction axiom. Include these \(J(c)\) and normalizations of the fixed
elementary arithmetic axioms. Denote the resulting recursive axiom
predicate by \(\operatorname{Base}(z)\). The normalization comes with
effective logical derivations in both directions, so this is an ordinary
presentation of \(\mathsf I\Sigma_1\).

Let \(E(e,z,t)\) be a primitive recursive computation-witness relation
for the fixed enumeration, so \(z\in W_e\) abbreviates
\(\exists t\,E(e,z,t)\). Define \[
U_e=\mathsf I\Sigma_1+
\{\sigma\in\mathcal P_3:\ulcorner\sigma\urcorner\in W_e\}.
\tag{12}
\] An extra-axiom line in a proof carries its accepting enumeration
computation; a base-axiom line carries its base certificate. Write
\(\operatorname{Prf}_e(p,z)\) for the resulting primitive recursive
proof-checking relation. Its existential provability predicate is
\(\Sigma_1\). This is a witnessed, or smooth, presentation. It also
satisfies the usual derivability conditions over \(\mathsf I\Sigma_1\).

Write \(\operatorname{RFN}_3(U_e)\) for \[
\forall z\forall p\,
\bigl((z\in\mathcal P_3\wedge\operatorname{Prf}_e(p,z))
\rightarrow\tau(z)\bigr).
\tag{13}
\] It expresses uniform reflection for the designated \(\Pi_3\)
sentences, not merely consistency. Choose a fixed-template prenex
normalization \(R_e\in\mathcal P_3\) of (13). The normalization is
provably equivalent in \(\mathsf I\Sigma_1\). A \(\Sigma_1\)
presentation of the proof predicate, negated in the antecedent, adds no
alternation beyond \(\Pi_3\).

The source theorem {[}PW{]} uses a language with exponentiation.
Appendix A verifies the conservative, complexity-preserving translation
between that convention and this one. We do not identify the two
signatures literally.

\subsection{Lemma 3. Uniform truth of the fixed
base}\label{lemma-3.-uniform-truth-of-the-fixed-base}

The partial truth predicate can be chosen so that \[
\mathsf{PA}\vdash
\forall z\bigl(\operatorname{Base}(z)\rightarrow\tau(z)\bigr).
\tag{14}
\] Moreover, \(\mathsf{PA}\) proves the compositional and substitution
equations for bounded-level satisfaction on subformulas of prenex
\(\Sigma_3\) and \(\Pi_3\) formulas.

\textbf{Proof.} Uniform evaluation of terms and bounded formulas has
matching existential and universal arithmetical descriptions. Arithmetic
proves their equivalence and totality. One obtains the fixed-level
satisfaction predicates by quantifying over the values of each of the at
most three unbounded blocks. The existential or universal description of
bounded satisfaction is chosen to match the last block. This gives the
required \(\Sigma_k\) and \(\Pi_k\) predicates, \(1\leq k\leq3\),
without an extra alternation. Arithmetic induction on the finite syntax
proves the term-substitution and compositional equations. The predicates
may be combined into a single tagged satisfaction formula for the union
of these classes; that combined formula need not itself be \(\Pi_3\).

Let \(S(c,n,a)=\operatorname{Sat}_{\Sigma_1}(c,n,a)\) express
satisfaction of the \(\Sigma_1\) formula coded by \(c\), with induction
variable \(n\) and remaining assignment \(a\). One instance of full
\(\mathsf{PA}\) induction, retaining \(c,a\) as parameters, proves \[
\forall c\forall a\Bigl(
\bigl[S(c,0,a)\wedge\forall n(S(c,n,a)\rightarrow S(c,Sn,a))\bigr]
\rightarrow\forall n S(c,n,a)\Bigr).
\tag{15}
\] The formula is guarded by the recursive condition that \(c\) codes a
suitable \(\Sigma_1\) formula. The uniform truth equations identify
(15), on such codes, with truth of \(J(c)\).

For completeness, the required complexity bound on \(J(c)\) follows
directly from writing an induction axiom for \(\phi\in\Sigma_1\) as \[
\neg\phi(0)\ \vee\
\exists n\bigl(\phi(n)\wedge\neg\phi(Sn)\bigr)
\ \vee\ \forall n\phi(n).
\tag{16}
\] The three disjuncts are respectively in \(\Pi_1\), \(\Sigma_2\) and
\(\Pi_2\). They admit a uniform \(\Pi_3\) prenex normalization,
including universal closure of the parameters. This normalizes a fixed
induction template, not an arbitrary formula into a falsely bounded
complexity class. The truth equivalence is verified with the same
fixed-level satisfaction equations. The remaining base axioms are fixed
arithmetic theorems. This proves (14) and the lemma. \hfill\(\square\)

\subsection{Lemma 4. From column truth to
reflection}\label{lemma-4.-from-column-truth-to-reflection}

Uniformly in \(e\), \[
\mathsf{PA}\vdash
\left[\forall z(z\in W_e\rightarrow\tau(z))
\rightarrow\operatorname{RFN}_3(U_e)\right].
\tag{17}
\]

\textbf{Proof.} Work in \(\mathsf{PA}\), assuming the antecedent. Let
\(p\) be a witnessed \(U_e\)-proof of a designated \(\Pi_3\) sentence
\(\sigma\). Collect the finite list \(\gamma_1,\ldots,\gamma_m\) of
nonlogical axioms used by \(p\), retaining their enumeration or base
certificates. Every \(\gamma_i\) is in \(\mathcal P_3\).

The effective deduction transformation yields a purely logical sequent
derivation of \[
\neg\gamma_1,\ldots,\neg\gamma_m,\sigma.
\tag{18}
\] Ordinary first-order cut elimination is formalizable in
\(\mathsf{EA}^+\), arithmetic with total superexponentiation {[}PW,
Section 2.1{]}. Therefore \(\mathsf{PA}\) verifies the transformation to
a cut-free derivation. This eliminates logical cuts; the induction
axioms were retained as nonlogical premises, not eliminated by a claim
about arithmetic proof strength.

The subformula property puts the formulas of the resulting derivation
among term instances of subformulas of the displayed \(\Sigma_3\) and
\(\Pi_3\) formulas. Use the tagged partial satisfaction predicate of
Lemma 3. Arithmetic induction on a cut-free proof verifies soundness of
its sequents: an atomic initial sequent is true, propositional rules
preserve the disjunction, existential rules use the value of their
witness term, and universal rules use arbitrary values for the fresh
eigenvariable. Equality rules are justified by arithmetic equality and
the substitution equations. The induction can be carried out uniformly
over the proof code and assignments, because the satisfaction formula
has fixed finite complexity.

Under the antecedent of (17), every extra \(\gamma_i\) is true by its
enumeration certificate; every base \(\gamma_i\) is true by Lemma 3.
Thus (18) yields \(\tau(\ulcorner\sigma\urcorner)\). The choices of
\(p\) and \(\sigma\) were arbitrary, giving (13).

Keeping the negative premises as separate sequent entries is important.
We have not called an arbitrary Boolean closure of all the axioms a
\(\Pi_3\) formula. The proof verifies a fixed bounded-complexity base
plus the specified column, and does not assert reflection for
unrestricted \(\mathsf{PA}\). \hfill\(\square\)

\subsection{Lemma 5. The arithmetic information retained by a
column}\label{lemma-5.-the-arithmetic-information-retained-by-a-column}

Suppose \(T\) contains \(\mathsf{PA}\), is sound in the standard
canonical structure, and \[
W_e=\{m:T\vdash\tau(\bar m)\}.
\tag{19}
\] Then \(U_e\) is standard-sound and has exactly the same designated
\(\Pi_3\) consequences as \(T\).

\textbf{Proof.} For a fixed designated \(\Pi_3\) sentence \(\sigma\),
(11) gives \[
T\vdash\sigma
\quad\Longleftrightarrow\quad
T\vdash\tau(\overline{\ulcorner\sigma\urcorner})
\quad\Longleftrightarrow\quad
\ulcorner\sigma\urcorner\in W_e.
\tag{20}
\] Thus every such \(T\)-consequence is literally an extra axiom of
\(U_e\), with an actual enumeration witness. Conversely each extra axiom
of \(U_e\) is a \(T\)-consequence; so are its \(\mathsf I\Sigma_1\) base
axioms. An ordinary arithmetic deduction from them is a valid deduction
in the original consequence relation. Hence every arithmetic
\(U_e\)-consequence is a \(T\)-consequence. Its axioms are true by
soundness of \(T\), so \(U_e\) is standard-sound.

This comparison is an external assertion about actual sentences and
proofs. It does not assert that \(U_e\) knows an index for \(T\), or
that either theory proves the displayed equality of consequence sets.
\hfill\(\square\)

\section{5. The uniform presentation and its reflection
sentences}\label{the-uniform-presentation-and-its-reflection-sentences}

Apply Lemma 2 to the fixed formula \(\tau\), and write \[
e_j=B(j,\ulcorner\tau\urcorner),\qquad j\geq1.
\tag{21}
\] The map \(j\mapsto e_j\) is recursive under the hypothesized chain,
but no arithmetic proof of its totality has been obtained. We therefore
use a program whose code can be compiled without waiting for this
computation to finish.

Let \(w(j)\) be the index of the following program:

\begin{quote}
Compute \(B(j,\ulcorner\tau\urcorner)\). If it returns \(e\), enumerate
\(W_e\).
\end{quote}

The fixed root index \(r\) and \(j\) are data in the program. Standard
parameter compilation makes \(j\mapsto w(j)\) primitive recursive
independently of whether the program runs to completion. Define \[
V_j=U_{w(j)}.
\tag{22}
\] These theories have a uniformly recursively enumerable axiom
presentation. Extensionally, \(V_j\) and \(U_{e_j}\) have the same
axioms, but that fact alone is not our justification for exchanging
their reflection sentences.

\subsection{Lemma 6. Pointwise provable comparison of the two
presentations}\label{lemma-6.-pointwise-provable-comparison-of-the-two-presentations}

For each fixed standard \(j\geq1\), \(\mathsf I\Sigma_1\) proves
equivalence of the witnessed provability predicates of \(V_j\) and
\(U_{e_j}\), preserving conclusion codes. Consequently it proves \[
\operatorname{RFN}_3(V_j)
\leftrightarrow\operatorname{RFN}_3(U_{e_j}).
\tag{23}
\]

\textbf{Proof.} Lemma 2 gives an actual finite terminating computation
of \(B(j,\ulcorner\tau\urcorner)\) with output \(e_j\). For this fixed
\(j\), arithmetic verifies that computation. Determinism proves
uniqueness of its output.

An accepting computation of the wrapper program consists of that initial
computation followed by an accepting computation for \(W_{e_j}\).
Conversely an accepting computation for \(W_{e_j}\) can be prefixed by
the verified initial computation to obtain one for the wrapper. These
are effective transformations of finite computation witnesses, verified
in \(\mathsf I\Sigma_1\). They prove \[
\forall z\bigl(z\in W_{w(j)}\leftrightarrow z\in W_{e_j}\bigr)
\tag{24}
\] there, together with the explicit witness transformations.

Apply these transformations to the extra-axiom certificates on every
line of a finite proof, leaving its arithmetic formulas, logical
inferences and base certificates unchanged. Induction on the proof
verifies both translations and preserves the conclusion. This proves
equivalence of the provability predicates. Reflection for either
presentation says that all its provable designated \(\Pi_3\) sentences
satisfy the same predicate \(\tau\). The provability equivalence
therefore gives (23).

The verification is pointwise: a different fixed finite computation may
be used for each \(j\). It does not presuppose that
\(\mathsf I\Sigma_1\) or \(\mathsf{PA}\) proves the uniform termination
of \(B\). \hfill\(\square\)

\section{6. The published arithmetic
obstruction}\label{the-published-arithmetic-obstruction}

\subsection{Theorem 7. Pakhomov--Walsh
non-descent}\label{theorem-7.-pakhomovwalsh-non-descent}

There is no recursively enumerable sequence \((S_n)_{n<\omega}\) of
recursively enumerable extensions of \(\mathsf{EA}\) such that \(S_0\)
is \(\Pi_3\)-sound and \[
S_n\vdash\operatorname{RFN}_{\Pi_3}(S_{n+1})
\qquad\text{for every }n.
\tag{25}
\]

This is {[}PW, Theorem 3.4{]}, with that paper's arithmetized proof
conventions. We invoke the theorem, rather than claiming a new
arithmetic non-descent result. Its hypotheses do not require
\(\mathsf{EA}\) to prove that the sequence descends, or require the
theories themselves to prove a total index-selection function.

We shall use the special case of uniformly recursively enumerable,
standard-sound extensions of \(\mathsf I\Sigma_1\). Appendix A explains
the arithmetic language translation and the agreement of the reflection
presentations needed to place our theories within this special case. The
stronger base also avoids a hidden collection assumption when witnessing
a finite proof's recursively enumerable axioms.

\section{7. Proof of Theorem 1 and the dissertation
corollary}\label{proof-of-theorem-1-and-the-dissertation-corollary}

\textbf{Proof of Theorem 1.} Suppose a family satisfying the theorem
exists. Construct the indices \(e_j\), the primitive recursive wrapper
indices \(w(j)\), and the theories \(V_j\) as above. Lemmas 2 and 5 show
that all \(V_j\) are standard-sound and extend the normalized
\(\mathsf I\Sigma_1\) base.

Fix \(j\geq1\). The predecessor endorsement in Lemma 2, applied to
\(j+1\), gives an actual theorem \[
T_j\vdash\forall x\bigl(K_{j+1}\tau(x)
\leftrightarrow x\in W_{\bar e_{j+1}}\bigr).
\tag{26}
\] Together with the known successor truthfulness (2), arithmetic
deduction inside \(T_j\) yields \[
T_j\vdash\forall x\bigl(x\in W_{\bar e_{j+1}}
\rightarrow\tau(x)\bigr).
\tag{27}
\] Since \(T_j\) contains \(\mathsf{PA}\), Lemma 4 gives \[
T_j\vdash\operatorname{RFN}_3(U_{e_{j+1}}).
\tag{28}
\] Lemma 6, used at the particular index \(j+1\), now gives \[
T_j\vdash\operatorname{RFN}_3(V_{j+1}).
\tag{29}
\]

The remaining issue is where this reflection is proved. Equation (29) is
so far a theorem of the modal theory \(T_j\), whereas Theorem 7 requires
a proof in the arithmetic theory \(V_j\).

Let \(R_{w(j+1)}\) be the designated \(\Pi_3\) normalization of the
reflection sentence in (29). The normalization equivalence is provable
in the base, so \[
T_j\vdash R_{w(j+1)}.
\] By (11), \(T_j\) therefore proves
\(\tau(\overline{\ulcorner R_{w(j+1)}\urcorner})\). By Lemma 2 this
means \[
\ulcorner R_{w(j+1)}\urcorner\in W_{e_j}=W_{w(j)}.
\tag{30}
\] Thus \(R_{w(j+1)}\) is an \textbf{actual added axiom} of \(V_j\),
with a finite enumeration witness. Appending the base proof of its
equivalence with the reflection sentence gives an actual finite proof \[
V_j\vdash\operatorname{RFN}_3(V_{j+1}).
\tag{31}
\] This argument uses the external column comparison to identify a
particular axiom. It does not put an unproved assertion about the whole
modal theory into \(V_j\).

The tail \(V_1,V_2,\ldots\) is uniformly recursively enumerable by its
compiled presentation (22). Its theories are standard-sound and satisfy
(31). With the convention verified in Appendix A, it contradicts Theorem
7. This proves Theorem 1. \hfill\(\square\)

\subsection{Corollary 8. Negative answer to Open Problem
3.2.1}\label{corollary-8.-negative-answer-to-open-problem-3.2.1}

In Alexander's language and standard canonical semantics, there is no
family of knowing machines in which each \(K_i\) knows the truthfulness
and a particular code of \(K_{i+1}\), as well as the Peano axioms for
the full language.

\textbf{Proof.} The language has effective operator names. By Definition
1.2.17 and Theorem 1.2.6 of {[}A13{]}, the root's consequences are
recursively enumerable. The canonical clause is (3), and the Truth axiom
of self-doubting knowledge makes all consequences true in the same
standard structure. Knowing the Peano axioms for the full modal language
includes knowing ordinary arithmetic \(\mathsf{PA}\). Definition 3.1.5
gives (1), as explained in Section 2.2. Knowing successor truthfulness
for all formulas supplies, in particular, (2) for \(\tau\). All the
hypotheses of Theorem 1 hold, which is impossible. \hfill\(\square\)

\section{8. Scope and nearby
questions}\label{scope-and-nearby-questions}

\subsection{8.1. How much of the code schema was
used?}\label{how-much-of-the-code-schema-was-used}

For the contradiction, Lemma 2 only needs to select the \(\tau\)-column
indices. Its recursive calls use a concrete family of formulas. At the
request for the \(i\)-th machine on the way to machine \(j\), the
formula is \[
\operatorname{Code}_{i+1,
\operatorname{Code}_{i+2,\ldots,
\operatorname{Code}_{j,\tau}}},
\qquad 1\leq i<j,
\tag{32}
\] with the indicated finite nesting and the canonical free-variable
renamings. At \(i=j\) the formula is \(\tau\). Each expression is
finite; its modal depth grows with \(j\). Code knowledge on this
recursive family of templates suffices. We do not claim that this family
is minimal.

In particular, the proof cannot be transferred merely by assertion to a
code schema restricted to arithmetic formulas or to a fixed finite modal
depth. No construction or impossibility result for those variants is
supplied here.

\subsection{8.2. An effectively named sequence is a real
hypothesis}\label{an-effectively-named-sequence-is-a-real-hypothesis}

The algorithm computes the literal operator name needed at every
recursive call. An arbitrary infinite path in a recursively enumerable
relation need not be computable. Therefore a family over an ill-founded
relation does not automatically provide the computably specified
successor sequence required here.

This distinction is relevant to {[}A20, Theorem 79{]}, which gives true
self-referential families for some ill-founded recursively enumerable
partial orders when the ordinal requirement is removed. Its existence
statement does not supply a computable descending path satisfying the
present hypotheses. Relabelling a noncomputable path as \(0,1,\ldots\)
is not an effective change of syntax and need not preserve recursive
enumerability of the root. Our conclusion is not a well-foundedness
theorem for every ordinal-free endorsement relation.

\subsection{8.3. What is imported and what is proved
here?}\label{what-is-imported-and-what-is-proved-here}

The arithmetic non-descent theorem is Pakhomov--Walsh's. Alexander
supplies the question, modal consequence framework and canonical
semantics. The local reduction consists of the root-based selection
procedure with its retained endorsements, the specified
arithmetic-column conversion, and the proof-presentation comparison
needed to apply non-descent. These are written out rather than hidden
behind references to earlier research notes.

The proof uses ordinary published facts about arithmetized syntax,
partial truth, logical cut elimination and incompleteness. It is a hand
proof; it does not inherit a formal certificate from the separate
repository developments concerning Questions 3-1-1 through 3-1-3. No
independent human mathematical review or determination of novelty is
asserted. The supplemental note records further consequences, with their
proofs and exact assumptions, in the same repository.

\section{Appendix A. Arithmetic language and reflection
conventions}\label{appendix-a.-arithmetic-language-and-reflection-conventions}

This appendix supplies the convention check used in Sections 4 and 6. It
is not a new hypothesis on the machines.

\subsection{Lemma A.1. Conservative language extension for the
application of
non-descent}\label{lemma-a.1.-conservative-language-extension-for-the-application-of-non-descent}

The special case of Theorem 7 used here applies to the uniformly
presented theories \(U_e\) in the language \(L_0\), with reflection as
in (13).

\textbf{Proof.} The language of \(\mathsf{EA}\) in {[}PW{]} includes
exponentiation. The theory \(\mathsf I\Sigma_1\) in \(L_0\) proves
totality and uniqueness of its usual \(\Sigma_1\) graph, so it admits a
conservative definitional extension by that function and the elementary
notational symbols in the source language. The definitions and their
proof translations are fixed once for every \(U_e\).

The complexity of the translated reflection class must also be checked.
Over \(\mathsf I\Sigma_1\), a bounded formula in the extended language
has equivalent \(\Sigma_1\) and \(\Pi_1\) translations into \(L_0\). To
construct them, give graph values to the added-function terms; use
uniqueness to obtain the universal alternative; and compose the
descriptions under bounded quantification. Finite collections of
evaluation witnesses are available in \(\mathsf I\Sigma_1\). The
constructions and their equivalences are uniform in bounded formula
codes. These are the usual \(\Delta_1\) definitional translations.

For a prenex extended-language \(\Pi_3\) sentence, choose the \(\Pi_1\)
version of its translated bounded matrix. Its final universal block
merges with the last universal block of the given sentence. The
translated sentence therefore remains \(\Pi_3\). In the dual
\(\Sigma_3\) case, use the \(\Sigma_1\) matrix translation. These
fixed-class translations preserve the corresponding partial truth
predicates, by induction through their three unbounded blocks. The
reverse embedding of \(L_0\) into the extended language is immediate.

Translate proofs using the graph definitions and their totality and
uniqueness proofs. An induction axiom for an extended-language
\(\Sigma_1\) formula translates to an induction instance for its
equivalent \(L_0\) \(\Sigma_1\) formula. Extra axioms are the same
\(\mathcal P_3\) sentences and retain their enumeration witnesses. Thus
the proof transformations are uniform in \(e\), preserve conclusions up
to the stated translations, and are verified in the fixed base. Together
with the partial-truth comparison, they transfer uniform \(\Pi_3\)
reflection between the two presentations. Ordinary proof predicates and
witnessed proof predicates likewise agree in this base: witnesses for
the finite list of enumerated axioms can be collected there. We never
use an arbitrary extensionally equivalent proof predicate.

Consequently a uniformly enumerable, sound sequence of \(L_0\) theories
as in (31) would yield one of sound extensions of \(\mathsf{EA}\) with
the reflection relation of Theorem 7. The published obstruction applies.
\hfill\(\square\)

\section{References}\label{references}

\textbf{{[}A13{]}} Samuel Allen Alexander. \emph{The Theory of Several
Knowing Machines}. PhD dissertation, The Ohio State University, 2013.
OhioLINK accession \textbf{osu1366362999}.
\href{https://rave.ohiolink.edu/etdc/view?acc_num=osu1366362999}{Institutional
record}. Main source locations: Theorem 1.2.6, printed pp.~9--10; Lemma
1.2.12 and Definitions 1.2.15--1.2.19, pp.~12--16; Definitions
3.1.4--3.1.5, p.~86; Theorem 3.1.6, pp.~87--88; Open Problem 3.2.1,
p.~89.

\textbf{{[}A20{]}} Samuel Allen Alexander. Self-referential theories.
\emph{Journal of Symbolic Logic} \textbf{85}(4) (2020), 1687--1716.
\href{https://arxiv.org/abs/2008.11535v1}{arXiv:2008.11535v1}, 11 August
2020. The comparison in Section 8 uses Corollary 78 and Theorem 79,
author-version pp.~23--25.

\textbf{{[}C00{]}} Timothy J. Carlson. Knowledge, machines, and the
consistency of Reinhardt's strong mechanistic thesis. \emph{Annals of
Pure and Applied Logic} \textbf{105} (2000), 51--82. Cited as a
predecessor of the modal framework, not as the source of the answer to
Open Problem 3.2.1.

\textbf{{[}PW{]}} Fedor Pakhomov and James Walsh. Reflection ranks and
ordinal analysis. \emph{Journal of Symbolic Logic} \textbf{86} (2021),
1350--1384.
\href{https://doi.org/10.1017/jsl.2020.9}{DOI:10.1017/jsl.2020.9}.
Inspected version:
\href{https://arxiv.org/abs/1805.02095v2}{arXiv:1805.02095v2}, 19
October 2019. The proof uses Section 2.1 for arithmetic/proof
conventions and Theorem 3.4, author-version pp.~10--11, for arithmetic
non-descent.

\end{document}
